% !TITLE 四块玉·别情
% !AUTHOR 关汉卿
% !DYNASTY 元
% !GENRE 散曲
% !ORDER 4
% !EDITION 《全元散曲》，隋树森编，中华书局（1964年版为候选）
% !EDITION_STATUS pending

\begin{poem}
自送别，心难舍，一点相思几时绝？
凭阑袖拂杨花雪\textsuperscript{1}。
溪又斜\textsuperscript{2}，山又遮\textsuperscript{3}，人去也！
\end{poem}

\begin{annotations}
\notetext{1}{杨花雪：杨花（柳絮）飘飞如雪花。“凭阑袖拂”——靠着栏杆，用袖子拂开飘来的杨花。}
\notetext{2}{溪又斜：溪水弯弯曲曲流向远方。斜（xiá），此处与“遮”“也”押韵。溪水蜿蜒，视线随之延伸。}
\notetext{3}{山又遮：山峦层层遮挡。溪水拐弯后，又被山挡住——远行的人已经消失在山水之外。}
\end{annotations}

\begin{appreciation}
写《窦娥冤》的关汉卿，是元曲四大家之首，笔下最擅长的是泼天的冤情与市井的泼辣；可他留下的小令里，却有《四块玉》这样一支清清淡淡的情词。【四块玉】是南吕宫的曲牌，别情才是曲题。

自送别，心难舍，一点相思几时绝？开口就是送别之后的事——人已经走了，曲子从自送别起笔，意思是前面的场面全略过，留下的只有散场之后的空。一点相思四个字要细看：相思明明缠得人透不过气她偏说只有一点。说得越轻，透出来的分量越重，这一句的小，正衬出底下那个几时绝的大——这样一点东西缠起来却偏偏没完没了。

凭阑袖拂杨花雪，是全曲唯一的一个动作。她靠着栏杆柳絮飘得像下雪深一层浅一层地扑过来，她抬起袖子，轻轻拂开。杨花拂去了还会再来，相思也是这样。一个寻常的挥手把满心说不出口的话都装了进去。杨花比雪，是有来历的：晋代的谢道韫答她叔父问雪，说未若柳絮因风起，从此柳絮与雪就结了缘。关汉卿写柳絮如雪，不写它好看，写它拂不尽。

末三句最见本色：溪又斜，山又遮，人去也。溪水拐一个弯，山又挡上来，人的身影就这样一截一截断掉，最后什么都没有了。又字用了两次：拐一道弯不够还要再拐一道；遮一座山不够，还要再遮一座——每一道弯每一座山都是送的人心里往下一沉。末了三个字人去也，是全曲最痛快也最难受的一声：口语，直说，没有半点修饰，好像一下把话说尽了也一下把心收住了。

我们在两汉卷读过《行行重行行》，汉代的女子说“弃捐勿复道，努力加餐饭”，千回百转的心事出口只剩一句劝他多吃一口饭。《四块玉》走的是另一路：元曲是唱给市井听的讲的是当场就要唱明白，所以它敢把话说到头，敢用一个也字把惆怅钉在那里。汉诗的相思收在肚子里元曲的相思摆在脸上；一个用加餐饭体面地收住，一个用人去也干脆地截断。雅和俗，隔着一千年，底下原是同一颗心。
\end{appreciation}
